信号平均是激光超声波检测中提高信噪比的事实标准方法，但其对检测吞吐量和材料完整性施加了根本限制。
由于激光产生的超声波本质上是低能量的，检测到的信号通常被背景噪声主导，信噪比仅随平均次数的平方根提高\cite{liu2022deep}。例如，将平均次数从256翻倍至1024，仅获得约3 dB的SNR增益，同时 acquisition time翻倍，且照射在试样表面的累积激光能量也翻倍。
这造成了一个严重的权衡：实现可靠缺陷检测所需的SNR水平——实践中通常需要256至1024次平均——需要使实时或高吞吐量筛查不切实际的检测时间，
以及可能导致敏感部件表面微烧蚀或涂层损坏的激光能量水平\cite{yang2023unsupervised}。问题在原位检测场景中进一步复杂化：此类场景 access受限，组件表面可能脆弱或已加工，重复激光照射要么物理上不可取，要么合同上被禁止。

传统去噪方法提供的缓解有限。维纳滤波和基于小波的阈值等线性滤波器可以抑制噪声，但频繁扭曲缺陷检测所依赖的诊断相关声学特征——到达时间、峰值幅度和波形形态\cite{chen2019wavelet}。BM3D等更先进的基于块的方法在图像域实现强噪声抑制，但在时域中难以保留超声波波形的稀疏、瞬态特性\cite{dabov2007bm3d}。现有深度学习方法显示出一般去噪的前景\cite{zhang2017denoising}，但通常使用大量平均参考信号的数据进行训练和评估，在应用于实际检测中遇到的低平均条件时产生性能差距\cite{liu2022deep}。结果是存在一个关键空白：需要高吞吐量的检测场景、涉及敏感或涂层表面的场景，或需要原位评估的场景无法依赖大量平均，
但当前去噪替代方案无法在保留缺陷表征所必需的声学特征的同时提供必要的SNR改善。
因此，迫切需要能够在大幅减少平均条件下提供可靠信号质量的去噪方法，使激光超声波检测能够避免传统信号平均的表面损伤和吞吐量损失。